\section{服务账户 (Service Accounts)}\label{sec:accounts}

正如我们已经指出的，JAM 中的服务在某种程度上类似于以太坊中的智能合约，因为它除了其他项目外还包括代码组件、存储组件和余额。与以太坊不同的是，代码分为两个隔离的入口点，各自具有自己的环境条件；一个是\emph{精炼}（Refinement），本质上是无状态的，发生在核心内；另一个是\emph{累积}（Accumulation），是有状态的，发生在链上。我们现在将关注后者。

服务账户保存在状态中的 $\accounts$ 下，是从服务标识符 $\serviceid$ 到指定与 JAM 协议相关的服务属性的命名元素元组的部分映射。形式化地：
\begin{align}\label{eq:serviceaccounts}
  \serviceid &\equiv \Nbits{32} \\
  \accounts &\in \dictionary{\serviceid}{\serviceaccount}
\end{align}

服务账户定义为存储字典 $\sa¬storage$、预映像查找字典 $\sa¬preimages$ 和 $\sa¬requests$、代码哈希 $\sa¬codehash$、余额 $\sa¬balance$ 和免费存储偏移量 $\sa¬gratis$ 的元组，以及两个代码 gas 限制 $\sa¬minaccgas$ 和 $\sa¬minmemogas$。我们还记录有关账户的某些使用特征：创建时的时隙 $\sa¬created$、最近一次累积的时隙 $\sa¬lastacc$ 和父服务 $\sa¬parent$。形式化地：
\begin{align}\label{eq:serviceaccount}
  \serviceaccount \equiv \tuple{\ \begin{aligned}
    \sa¬storage &\in \dictionary{\blob}{\blob}\,,\
    \sa¬preimages \in \dictionary{\hash}{\blob}\,,\\
    \sa¬requests &\in \dictionary{\tuple{\hash,\bloblength}}{\sequence[:3]{\timeslot}}\,,\\
    \sa¬gratis &\in \balance\,,\
    \sa¬codehash \in \hash\,,\
    \sa¬balance \in \balance\,,\
    \sa¬minaccgas \in \gas\,,\\
    \sa¬minmemogas &\in \gas\,,\
    \sa¬created \in \timeslot\,,\
    \sa¬lastacc \in \timeslot\,,\
    \sa¬parent \in \serviceid\\
    %i, o, f
  \end{aligned}\,}
\end{align}

因此，索引为 $s$ 的服务的余额记为 $\accounts\subb{s}_\sa¬balance$，该服务键 $\mathbf{k} \in \blob$ 的存储项记为 $\accounts\subb{s}_\sa¬storage\subb{\mathbf{k}}$。




\subsection{代码与 Gas (Code and Gas)}

服务账户的代码和相关元数据由哈希标识，如果该服务要正常运作，该哈希必须存在于其预映像查找中（参见第 \ref{sec:lookups} 节），并且预映像是两个数据块的正确编码。因此我们定义实际代码 $\sa¬code$ 和元数据 $\sa¬metadata$：
\begin{align}
  \forall \mathbf{a} \in \serviceaccount : \tup{\mathbf{a}_\sa¬metadata, \mathbf{a}_\sa¬code} \equiv \begin{cases}
    \tup{\mathbf{m}, \mathbf{c}} &\when \encode{\var{\mathbf{m}}, \mathbf{c}} = \mathbf{a}_\sa¬preimages[\mathbf{a}_\sa¬codehash] \\
    \tup{\none, \none} &\otherwise
  \end{cases}
\end{align}

代码中有两个入口点：
\begin{description}
  \item[0 \texttt{refine}]精炼，在核心内执行且无状态。\footnote{技术上有一些小的状态假设，即每个服务的预映像的适度最新实例。具体细节在第 \ref{sec:packagesanditems} 节中讨论。}
  \item[1 \texttt{accumulate}] 累积，在链上执行且有状态。
\end{description}

精炼和累积分别在第 \ref{sec:computeworkreport} 和 \ref{sec:accumulationexecution} 节中有更详细的描述。

如附录 \ref{sec:virtualmachine} 所述，JAM 虚拟机中的执行时间以\emph{gas}为单位确定性地测量，表示为小于 $2^{64}$ 的自然数，形式化地记为 $\gas$。如果数量可以为负，我们也可以使用 $\signedgas$ 表示集合 $\Z_{-2^{63}\dots2^{63}}$。账户中指定了两个限制，它们确定执行服务代码的\emph{累积}入口点所需的最小 gas。$\sa¬minaccgas$ 是每个工作项所需的最小 gas，而 $\sa¬minmemogas$ 是每个延迟传输所需的最小 gas。




\subsection{预映像查找 (Preimage Lookups)}\label{sec:lookups}

除了在仅链上可用的任意键/值对中存储数据之外，账户还可以请求数据在核心内也可用，从而可用于服务代码的精炼逻辑。关于此功能的状态保存在服务的 $\sa¬preimages$ 和 $\sa¬requests$ 组件下。

预映像查找和存储之间有几个区别。首先，预映像查找充当从哈希到其预映像的映射，而通用存储将任意键映射到值。其次，预映像数据是外部提供的，而存储数据源自服务的累积。第三，预映像数据一旦提供，不能自由删除；它经历一个被标记为不可用的过程，只有在一段时间后才能从状态中移除。这确保了有关其存在的历史信息得以保留。最后一点尤其重要，因为预映像数据被设计为在核心内查询，在服务代码的精炼逻辑下，因此预映像的历史可用性是已知的这一点非常重要。

我们首先重新制定有关数据查找系统的状态部分。该系统的目的是提供一种在链上存储静态数据的方式，使其随后可以作为仅接受数据哈希及其八位字节长度的函数在任何服务代码的执行中可用。

在\emph{累积}函数的链上执行期间，这是微不足道的，因为验证区块的所有验证者必然具有完整知识的内在状态，即 $\thestate$。然而，对于\emph{精炼}的核心内执行，没有所有验证者内在可用的此类状态；因此我们命名一个历史状态——\emph{查找锚点}——它必须在工作的影响可以累积之前被视为最近已最终确认，从而提供此保证。

通过保留有关其可用性的历史信息，我们确信任何具有链的最近最终确认视图的验证者能够确定任何给定预映像在审计可能发生的时期内任何时间是否可用。这确保了判决将具有确定性的信心，即使在链状态上没有共识。

重述一下，我们必须能够定义某个\emph{历史}查找函数 $\histlookup$，它确定某个哈希的预映像在某个时隙是否可供某个服务账户查找，如果是，则提供它：
\begin{equation}
\begin{aligned}
  \histlookup\colon \abracegroup[\ ]{
    \tuple{\serviceaccount, \Nmax{(\H_\¬timeslot - \Cexpungeperiod) \dots \H_\¬timeslot}, \hash} &\to \optional{\blob} \\
    (\mathbf{a}, t, \blake{\mathbf{p}}) &\mapsto v : v \in \set{ \mathbf{p}, \none }
  }
\end{aligned}
\end{equation}

该函数在下方公式 \ref{eq:historicallookup} 中定义。

索引为 $s$ 的服务的预映像查找记为 $\accounts\subb{s}_\sa¬preimages$，是一个将哈希映射到其对应预映像的字典。此外，与查找关联的元数据记为 $\accounts\subb{s}_\sa¬requests$，是一个将某个哈希和假设长度映射到历史信息的字典。

\subsubsection{不变量 (Invariants)}

查找系统的状态自然满足若干不变量。首先，任何预映像值必须对应其哈希，公式 \ref{eq:preimageconstraints}。其次，预映像值存在于状态中意味着其哈希和长度对有某种关联状态，也在公式 \ref{eq:preimageconstraints} 中。形式化地：
\begin{equation}\label{eq:preimageconstraints}
  \forall \mathbf{a} \in \serviceaccount, \kv{h}{\mathbf{d}} \in \mathbf{a}_\sa¬preimages \Rightarrow
    h = \blake{\mathbf{d}}\wedge
    \tup{h , \len{\mathbf{d}}} \in \keys{\mathbf{a}_\sa¬requests}
\end{equation}

\subsubsection{语义 (Semantics)}

历史状态组件 $h \in \sequence[:3]{\timeslot}$ 是最多三个时隙的序列，该序列的基数暗示四种模式之一：
\begin{itemize}
  \item{$h = \sequence{}$}：预映像已被\emph{请求}，但尚未提供。
  \item{$h \in \sequence[1]{\timeslot}$}：预映像\emph{可用}，自时间 $h_0$ 起。
  \item{$h \in \sequence[2]{\timeslot}$}：先前可用的预映像现在自时间 $h_1$ 起\emph{不可用}。它自时间 $h_0$ 起曾可用。
  \item{$h \in \sequence[3]{\timeslot}$}：预映像\emph{可用}，自时间 $h_2$ 起。它先前自时间 $h_0$ 至 $h_1$ 期间曾可用。
\end{itemize}

历史查找函数 $\histlookup$ 现在可以定义为：
\begin{equation}
  \begin{aligned}\label{eq:historicallookup}
    &\histlookup\colon \tuple{\serviceaccount, \timeslot, \hash} \to \optional{\blob} \\
    &\histlookup(\mathbf{a}, t, h) \equiv \begin{cases}
      \mathbf{a}_\sa¬preimages\subb{h}\!\!\! &\when h \in \keys{\mathbf{a}_\sa¬preimages} \wedge I(\mathbf{a}_\sa¬requests\subb{h, \len{\mathbf{a}_\sa¬preimages\subb{h}}}, t) \!\!\!\! \\
      \none &\otherwise
    \end{cases}\\
    &\where I(\mathbf{l}, t) = \begin{cases}
      \bot &\when \sq{} = \mathbf{l} \\
      x \le t &\when \sq{x} = \mathbf{l} \\
      x \le t < y &\when \sq{x, y} = \mathbf{l} \\
      x \le t < y \vee z \le t &\when \sq{x, y, z} = \mathbf{l} \\
    \end{cases}
  \end{aligned}
\end{equation}




\subsection{账户占用与阈值余额 (Account Footprint and Threshold Balance)}

我们定义依赖值 $\sa¬items$ 和 $\sa¬octets$ 为服务的存储占用，具体是存储中的项目数和存储中使用的总八位字节数。它们纯粹由服务的存储映射定义，并且必须假设每当服务的存储被更改时，这些也会改变。

此外，正如我们将在第 \ref{sec:serialization} 节的账户序列化函数中看到的，这些预计会显式地出现在 Merklized 状态数据中。因此我们明确地设定它们的集合。

然后我们可以定义第三个依赖术语 $\sa¬minbalance$——任何给定服务账户所需的最小或\emph{阈值}余额，根据其存储占用。
\begin{align}
  \forall \mathbf{a} \in \values{\accounts}\colon \abracegroup{
    \mathbf{a}_\sa¬items \in \Nbits{32} &\equiv
      2\cdot\len{\,\mathbf{a}_\sa¬requests\,} + \len{\,\mathbf{a}_\sa¬storage\,} \\
    \mathbf{a}_\sa¬octets \in \Nbits{64} &\equiv
      \sum\limits_{\,\tup{h, z} \in \keys{\mathbf{a}_\sa¬requests}\,} \!\!\!\!81 + z \\
    &\phantom{\equiv\ } + \sum\limits_{\tup{x, y} \in \mathbf{a}_\sa¬storage} 34 + \len{y} + \len{x} \\
    \label{eq:deposits}
    \mathbf{a}_\sa¬minbalance \in \balance &\equiv
      \max(0,
        \Cbasedeposit
        + \Citemdeposit \cdot \mathbf{a}_\sa¬items
        + \Cbytedeposit \cdot \mathbf{a}_\sa¬octets
        - \mathbf{a}_\sa¬gratis
      )
  }
\end{align}



\subsection{服务特权 (Service Privileges)}
JAM 包含向若干服务授予特权的能力。持有此信息的状态部分记为 $\privileges$，包含五种特权。第一个 $\manager$ 是\emph{管理器}服务的索引，该服务能够从区块到区块更改 $\privileges$ 以及向服务授予存储存款信用。下一个 $\delegator$ 能够设置 $\stagingset$。然后 $\registrar$ 单独能够在受保护范围内创建新服务账户索引。接下来 $\assigners$ 是能够更改授权器队列 $\authqueue$ 的服务索引，每个核心一个。

最后，$\alwaysaccers$ 是一个小字典，包含在每个区块中自动累积的服务索引以及每个服务累积使用的基本 gas 量。形式化地：
\begin{align}
  \label{eq:privilegesspec}
  \privileges &\equiv \tuple{
    \manager,
    \delegator,
    \registrar,
    \assigners,
    \alwaysaccers
  }\\
  \manager &\in \serviceid \ ,\qquad
  \delegator \in \serviceid \ ,\qquad
  \registrar \in \serviceid \\
  \assigners &\in \sequence[\Ccorecount]{\serviceid} \ ,\qquad
  \alwaysaccers \in \dictionary{\serviceid}{\gas}
\end{align}
